\section{特征提取}

语音信号经过预处理和短时分析后，需要进一步提取能够表征不同数字之间声学差异的特征。
梅尔频率倒谱系数（Mel-Frequency Cepstral Coefficients, MFCC）是语音识别中最广泛应用的特征之一，
其设计模拟了人耳对声音的非线性感知特性。
本节详细介绍 MFCC 特征提取的完整流程，包括预加重、分帧加窗、FFT 与功率谱、梅尔滤波、对数压缩与离散余弦变换，以及动态特征的构建。

\subsection{预加重}

语音信号在唇辐射和空气传输过程中，高频分量会自然衰减。
为补偿这一效应，首先对离散语音信号 $x[n]$ 进行预加重（pre-emphasis），即通过一阶 FIR 高通滤波器：
\begin{equation}
  y[n] = x[n] - \alpha\,x[n-1],
\end{equation}
其中预加重系数 $\alpha = 0.97$。
该滤波器的作用是增强信号的高频成分，使频谱在高频段更加平坦，有利于后续的频谱分析。
从信号与系统的角度看，该系统的传递函数为 $H(z) = 1 - \alpha z^{-1}$，是一个简单的一阶高通 FIR 滤波器。

\subsection{分帧、加窗与短时傅里叶变换}

语音信号是典型的非平稳信号，其统计特性随时间变化。
但在较短的时段（约 20--30 ms）内，发音器官的运动相对缓慢，信号可近似视为平稳。
基于这一短时平稳假设，对预加重后的信号进行分帧处理：采用窗长为 25 ms 的 Hamming 窗，帧移为 10 ms。
在 16 kHz 采样率下，窗长对应 $N = 400$ 个采样点，帧移对应 $160$ 个采样点。

Hamming 窗的时域表达式为
\begin{equation}
  w[n] = 0.54 - 0.46\cos\left(\frac{2\pi n}{N-1}\right),\quad 0 \le n \le N-1.
\end{equation}
加窗可以减小帧边界处的不连续性，抑制频谱泄漏。

对每一帧加窗后的信号进行 $K$ 点 FFT，得到短时幅度谱 $|X(m,k)|$，进而计算功率谱 $|X(m,k)|^2$。
关于 STFT 的详细讨论见语音信号分析一节。

\subsection{梅尔滤波器组}

人耳对声音频率的感知是非线性的：在低频区域，人耳能够分辨较细微的频率差异；而在高频区域，频率分辨率则相对较低。
梅尔（Mel）尺度正是为了模拟这一感知特性而提出的一种频率映射，其与物理频率 $f$（Hz）的近似关系为
\begin{equation}
  \mathrm{Mel}(f) = 2595\log_{10}\!\left(1 + \frac{f}{700}\right).
\end{equation}

在梅尔尺度上等间距地放置一组三角滤波器，再映射回线性频率轴，即构成梅尔滤波器组。
梅尔滤波器在低频段分布密集、在高频段分布稀疏，体现了人耳的非线性频率分辨率。

本实验采用 40 个三角梅尔滤波器，频率范围为 20 Hz 至 7600 Hz（对应 16 kHz 采样的奈奎斯特频率以下）。
每一帧的功率谱通过梅尔滤波器组后，得到 40 维的对数梅尔能量。

\subsection{对数压缩与离散余弦变换}

将梅尔滤波器组的输出取对数，一方面模拟人耳对响度的对数感知特性，另一方面压缩动态范围，使特征对输入能量的绝对大小不敏感：
\begin{equation}
  S(m, l) = \log\left(\sum_{k=0}^{K-1} |X(m,k)|^2 \cdot H_l[k]\right),
\end{equation}
其中 $H_l[k]$ 为第 $l$ 个梅尔滤波器的频率响应，$l = 0, 1, \dots, L-1$，$L = 40$。

对数梅尔能量各维度之间存在较强的相关性。为降低特征维度并去除相关性，对 $S(m, l)$ 沿滤波器轴进行离散余弦变换（DCT）：
\begin{equation}
  c(m, i) = \sum_{l=0}^{L-1} S(m, l) \cos\!\left[\frac{\pi i}{L}\left(l + \frac{1}{2}\right)\right],\quad i = 0, 1, \dots, N_{\mathrm{mfcc}}-1.
\end{equation}
DCT 是一种正交变换，能够将能量集中到少数低阶系数上，同时使各阶倒谱系数近似不相关。
保留前 $N_{\mathrm{mfcc}} = 13$ 个 DCT 系数，即为该帧的 MFCC 特征。

值得注意的是，第 0 阶系数 $c(m,0)$（常记为 C0）主要反映该帧的整体能量，数值通常远大于其余系数。
在后续特征可视化中，为避免 C0 掩盖其他倒谱系数的细节结构，13 维 MFCC 图在绘制时省略 C0，仅显示 C1--C12。

\begin{figure}[H]
  \centering
  \includegraphics[width=0.75\linewidth]{signal_mfcc.png}
  \caption{13 维 MFCC 特征图（可视化时省略 C0，仅显示 C1--C12）。横轴为帧序号，纵轴为 MFCC 系数阶数，颜色表示系数值的大小。}
  \label{fig:mfcc13}
\end{figure}

\subsection{动态特征}

静态 MFCC 仅反映当前帧的频谱包络信息，无法刻画语音信号随时间变化的动态特性。
为此，在 13 维静态 MFCC 的基础上，分别计算一阶差分（Delta, $\Delta$）和二阶差分（Delta-Delta, $\Delta^2$）：
\begin{equation}
  \Delta_t = \frac{\sum_{n=1}^{N_d} n\,(c_{t+n} - c_{t-n})}{2\sum_{n=1}^{N_d} n^2},
\end{equation}
其中 $c_t$ 为第 $t$ 帧的 MFCC 向量，$N_d$ 为差分窗口半径（本实验取 $N_d = 2$）。
$\Delta^2$ 以同样的方式在 $\Delta$ 序列上计算得到。

最终，每一帧的特征向量由 13 维静态 MFCC、13 维 $\Delta$ 和 13 维 $\Delta^2$ 拼接而成，共计 39 维帧级特征。
动态特征的引入使模型能够同时利用频谱包络的静态信息和帧间变化的动态信息，从而更好地建模语音的时序演变规律。

\begin{figure}[H]
  \centering
  \includegraphics[width=0.75\linewidth]{signal_mfcc_features.png}
  \caption{39 维 MFCC + $\Delta$ + $\Delta^2$ 特征图。上段为 13 维静态 MFCC（含 C0），中段为 13 维 $\Delta$，下段为 13 维 $\Delta^2$。横轴为帧序号，纵轴为系数阶数。}
  \label{fig:mfcc39}
\end{figure}

图 \ref{fig:mfcc13} 和图 \ref{fig:mfcc39} 分别展示了同一语音段在两种特征表示下的可视化结果。
13 维特征图（图 \ref{fig:mfcc13}）聚焦于静态频谱包络的时变模式，不同数字的 MFCC 在特定系数阶数上呈现出差异化的能量分布；
39 维特征图（图 \ref{fig:mfcc39}）则进一步整合了动态信息，$\Delta$ 和 $\Delta^2$ 分量反映了频谱变化的速率与加速度，
为后续的分类模型提供了更丰富的时序判别依据。

\subsection{与信号与系统课程的关联}

MFCC 特征提取流程与信号与系统课程的知识体系紧密相关。
预加重对应一阶 FIR 滤波器的设计与频率响应分析；
分帧加窗涉及信号截断与频谱泄漏问题，是离散时间信号处理的重要内容；
FFT 和功率谱估计是频域分析的核心工具，在语音信号分析一节中已有详细讨论；
梅尔滤波器组体现了频域加权的思想，本质上是对功率谱的频域滤波；
对数运算实现动态范围压缩，与信号的非线性变换相关；
DCT 作为正交变换，具有能量集中和去相关的优良性质，是信号表示理论中的典型方法。
这些信号处理技术的有机结合，构成了从原始语音波形到紧凑声学特征表示的关键桥梁。
